\section{Hercules’ Seventh Triumph}

As I have argued, Hercules has decisively ``come into his own'': he has achieved the unstated goal of the classical \emph{polis}, which was to become an \emph{Aristotelian unequal}. He is not yet equal to gods, but he no longer lives under mortal law, because he has suffered its stroke, and lived (which is the only way to transcend it). He has passed through the difficult time of his hero saga. In any heroic legend, there are moments, particularly early on, when the hero is set upon by all seemingly legitimate forces. In doing so, these forces either illegitimate themselves, or show themselves to have been (in fact) illegitimate from the beginning anyway: they ``fall away'' before the great individuality.

Traditionalists should never forget that King Arthur, King Alfred\footnote{\url{https://en.wikipedia.org/wiki/Alfred\_the\_Great}}, Robin Hood, and any other number of folk heroes, saints, and strong men were initially persecuted, not merely or even especially by the wicked, but by their ``own''. They are inevitably portrayed as a ``wolf's head'', painted as demons or villains, and an attempt is made to exterminate them in the name of all that is sacred. Often, this attempt actually succeeds, and gives us legends of martyrs (at best) or legends of villains (at worst). Thus, the entire nation of Germany (for example) occupies the dubious position of being an arch-villain in the modern narrative, in which Kaisers and Barons and tool mechanics from the Wehrmacht fare no better than Hitler himself. The modern world has spent enormous spiritual energy probing the ``dark side'' of the nobility of Old Europe.

How many movies\footnote{\url{https://www.youtube.com/watch?v=QpDEMKuzixc}} have given us the stereotypical aristocrat who is bent on murdering and pillaging and oppressing those unlucky enough to fall under his aegis? Until an equivalent effort is made to diagnose and bring to light the even more hideous brutalities of the mob during the modern era (does anyone even count the number of automobile accidents on the roads against the ``age of the Machine''?) we will not be remotely out of the woods. And there is a disadvantage --- the educated classical liberals were well-equipped spiritually and intellectually to critique their predecessors, whereas we are much less able to critique the enormities that have emerged in the modern age of the Mass.

So Traditionalists should expect to encounter massive reserves of energy that will attempt to blindly preserve their grip on power. This is one of the reasons I warned in one of my posts against violating General Law in an indiscriminate way: the mindset of an era represents a spiritual energy (however degenerate and blind) which is guaranteed to activate and suppress dissidence of any kind. We should be more adept at ``reading the times and seasons'', and knowing what can be accomplished, and how it should be done. A great many people will make the mistake of becoming Don Quixotes, and end up being battered by windmills. It will not be easy, quick, or safe to extricate mankind from the spiritual darkness it has fallen into, over two centuries or more, in the West. Hercules started out as an underdog, with the gods against him, the nobility of the time actively seeking his death, and the common people (no doubt) simply channeling their basic desire to ``witch-hunt'' into the approved mode\footnote{\url{http://unqualified-reservations.blogspot.com/2013/09/technology-communism-and-brown-scare.html}}. Interested readers might wish to consult the subject of Ponerology\footnote{\url{http://www.amazon.com/Political-Ponerology-Science-Adjusted-Purposes/dp/1897244258}} for a closer look at this dynamic. In any case, in order to be a hero, you have to be willing to be painted (if the need arises) as a villain, perhaps even by those closest to you, not for the sake of the black hat, but out of a clear apprehension of what is right. Read King David's Psalms, for instance, for more examples.

You must be willing to stand against the world, like Athanasius, if you wish to overcome it. Indeed, how could it be otherwise? In our time, the one position that is guaranteed to be a challenge to ``all-comers'' is the position of antique Europeanism, in the chivalric Christian tradition. It will infuriate almost anyone, for (indeed), we have all been very subtly programmed against it, given antidotes and immunization shots to prevent ``catching it'', and taught in the very bosoms of our Churches that such a world should have never existed (at best) or even did not exist (at worst). Go and stand in some of the older Churches\footnote{\url{http://www.urlaube.info/en/Wien/Kirchen.html}} of Europe or North America. The pictures of the monks or founders or patrons tell a very different story from the bleak world we are taught that they inhabited. Read an older and sympathetic novel or biography of that period -- something irrevocable has been publicly lost from that period. The clear look, the firm jaw, the clean heart -- they are gone from our consciousness. They are submerged.

It is to restore such that Hercules comes. Not to achieve immortality directly, but to right what is wrong, and to spill sweat, tears, toil, and blood. Will he come unopposed and welcomed as a hero? He comes tasked with labors, and against the ``odds and the gods'', against heaven and earth. Not everyone has to be a public Hercules, but everyone has a dragon or two in his circle that needs a-killin'. Rest assured, the dragon knows your name, and will exploit this. How else can evil exist, except to stay in the dark, and how else can evil have union, except in the dark? And who does not have dark within them? This will be an arduous progression.

For the seventh labor, Hercules is sent against the Cretan Bull. It is tearing up vineyards and terrorizing the island. When Martin Luther reared his head in the wolds of Germany, the pope noticed and said ``a wild boar is loose in God's vineyard''. We do not know all of what Luther had to reckon with, only that his coming marked the end of a unified Western world\footnote{\url{http://www.amazon.com/Three-Reformers-Luther-Descartes-Rousseau/dp/0804610800}}. When the Reformation was finished, the age of Atheism and Liberalism began\footnote{\url{http://www.righthandpath.org/the-church-and-the-common-mind/}}. It may seem harsh to say this, and there are no doubt many things to salvage from Northern European spirituality, but Luther left Europe in rags and flames.

Interestingly enough, Hercules submits the bull of Cretes in a strangle-hold, refusing to kill it, instead sending it back to his royal tormentor, who wishes to offer it to Hera as a sacrifice. Hera declines to do this, and is offended, so the bull is shipped off somewhere else, and ends up being sacrificed to Athena and Apollo. Hercules has restored the natural order of things, the Logos, since Athena (wisdom) and Apollo (war/arts) are higher divinities than the jealous Hera, who is obsessed eternally (much like a feminist) with ancient and perceived wrongs and injustices that can never possibly be ``righted'' except by making more suffering \emph{now}.

The bull is associated with animal-instincts and drives. Not only was Crete known in Biblical times as a lustful and lazy island, but the bull itself wreaked a moral havoc on the order of the island:

\begin{quotex}
Minos himself, in order to prove his claim to the throne, had promised the sea-god Poseidon that he would sacrifice whatever the god sent him from the sea. Poseidon sent a bull, but Minos thought it was too beautiful to kill, and so he sacrificed another bull. Poseidon was furious with Minos for breaking his promise. In his anger, he made the bull rampage all over Crete, and caused Minos' wife Pasiphae to fall in love with the animal. As a result, Pasiphae gave birth to the Minotaur, a monster with the head of a bull and the body of a man. Minos had to shut up this beast in the Labyrinth, a huge maze underneath the palace, and every year he fed it prisoners from Athens.
\end{quotex}


The broken promises and moral decay had birthed the Minotaur, a beast which would also have to be slain by another hero. But that is another story. And here is a goodly lesson --- it is better not to vow, than to vow, and not pay. If we fight with monsters, the only way to not become one is to keep one's vows. This relates not merely to not using the dark arts (for one is already pledged to what is noble and high) but also to wisely watching one's actions, not thinking more highly of Self than one ought. Some monsters are other people's property --- it is their destiny to slay them, and they are chosen to do so. If you wish to aid, aid them --- do not usurp their rightful place as slayer of the beast. Slay your own beasts. The restoration of a proper moral Order in the thoughts and hearts of those who oppose the modern temper and spirit is the \emph{sine qua non} of Victory, since the Victory comes from above, vertically, just as initiation does.

Becoming Captain Ahab or Don Quixote will not do anyone much good. We should be looking around for the chosen ones, all the while, doing all we can to make sure we can hear the still, small voice ourselves, for God will surely call each of us to some labor of Hercules, big or small. In short, the restoration of Order will come in a divine web or hierarchy of spiritual intervention, with the sacred hero-priest-king as the focal point. Squabbling and in-fighting are trademarks of the Left, what thieves and criminals do when they have despoiled someone and are dividing the loot. A man of the Right stands for the invisible and absolute order which he alone discerns, and suffers to make visible. The unveiling of this invisible web, through each one of us, will annihilate the plots of darkness, but it has to be (as the old preacher used to put it) ``all of grace'' or ``of God''. Here we see that Christianity still has a mission for the ``New Right'' -- until something greater than it is found (and how could that happen, without victory?), victory will require a re-harmonization between Christian Tradition and the remnants of the West. \emph{This will be resisted on all sides, and itself proves that it is the salient point.}

It is entirely possible that Hercules could have died or failed, in which case, he would be portrayed as a massive freak and fraud. This risk will have to be born up under, but born up under wisely, and strongly, and with the best possible chance of success. This is the path of the strong --- not to fight invincibly and surely, but to triumph over a multitude of doubts and anguish and sins within, and to win through (if God wills) against a greater multitude of such externally, to a doubtful and dark destiny, under the shade of the laurel and the oak and perhaps the cypress. Invincibility is coming, but for now, is not yet.

One of the greatest strengths and services and preparations for this is to read the exploits of the past and to ``right the wrong'' of the ancestors. Look at the Iron Guard\footnote{\url{https://en.wikipedia.org/wiki/Iron\_Guard}}, for example, and their fate. How could they have failed? Most men are simply unwilling to risk a similar fate, or even honor their memory, and so they conform in their thoughts. The man of the Right pays respect and honor to his ancestors, which means learning their stories and ``righting the wrong'' in his heart, both by avoiding their mistakes and faults, but also by re-baptizing their memory and burnishing it till it shines brightly. Only thus, in the path of the Past, can we expect to find blessing for the dangerous Future, and strength for the present. This will be unappealing to the immense majority of men, who are depressed by such efforts, and cling to a hope for a brighter future through democracy, chemistry, and modern organization. They are often more decent, more intelligent, and more accomplished than many who see the beauty of the Past for today. There are many emoluments, amenities, graces, and goods that accrue to those who forsake this fight, however reluctantly and with whatever good intentions. Others will mock the futility of the ``martyrs'', and use it as an occasion to desert the men of honor from the past.

Nevertheless, after all, Zeus will be vindicated, and He is not mocked. Who will wear his laurel leaves? He who struggles to the end.

There are compensations. After each labor, you get to see Eurystheus hiding in his \emph{piloi}, and hear Hera cackling with rankled and wounded pride. Zeus will send his compliments, secretly. Those who have suffered will acclaim you as their local hero. You will be allowed the noble and theatrical gesture of \emph{noblesse oblige}\footnote{\url{https://en.wikipedia.org/wiki/Noblesse\_oblige}}. After all, who asked them to anoint themselves as god of this world? Who asked them to invade the past and obliterate it?

\begin{quotex}
\begin{verse}
But the Consul's\footnote{\url{http://ancienthistory.about.com/library/bl/bl\_horatiuspoem.htm}} brow was sad, and the Consul's speech was low,\\
And darkly looked he at the wall, and darkly at the foe.\\
``Their van will be upon us before the bridge goes down;\\
And if they once might win the bridge, what hope to save the town?''

Then out spoke brave Horatius\footnote{\url{http://ancienthistory.about.com/od/romanfamousfigures/g/Horatius.htm}}, the Captain of the Gate:\\
``To every man upon this earth, death cometh soon or late;\\
And how can man die better than facing fearful odds,\\
For the ashes of his fathers, and the temples of his Gods,

And for the tender mother who dandled him to rest,\\
And for the wife who nurses his baby at her breast,\\
And for the holy maidens who feed the eternal flame,\\
To save them from false Sextus, that wrought the deed of shame?
\end{verse}
\end{quotex}


A man like that might make a difference, no matter what may come. That is why creating a spiritual \emph{elite} (nobility) is the first task of the Counter-Revolution. Those who wish for adventure, to see brave deeds and to truly live, will be drawn to their course, and will uncover the stone of their destiny, which may be dark or bright, victorious or defeated, but which will be uniquely their own, with their own name written upon it\footnote{\url{http://biblehub.com/revelation/2-17.htm}}. And this stone is part of the final victorious Temple, so that the death of a Byrhtnoth\footnote{\url{http://en.wikipedia.org/wiki/Battle\_of\_Maldon}} is still triumphant.

Let me sum it up to be crystal clear --- the Hero will fight wisely, not beating the air, with the best chance of success, but he will also fight nobly, not avoiding an inevitable or symbolic death if the situation calls for such. God is able to equally deliver him, in either way, either before or after his physical demise or failure.

Hercules' seventh triumph is a public act of wrestling the immorality and violence of an entire island to its knees, and sending the trophy as a taunt and a gift to his persecutors. The death of the bull, sacrificed to two war gods of wisdom and the fine arts, founds an era of moral stability, which the slayer of the Minotaur will lengthen and perpetuate. He is almost invincible, and is now acting almost openly on behalf of Zeus.

Could a determined alliance between a remnant Church and honorable pagans (who understand the divine mission of Christianity) wrestle the bull to its knees, in the tradition of Mithras, and of the Christ of Revelation? If enough men advance through the Labors, it will be a real possibility. Perhaps there is a man out there, reading this, who has within him both the Church and Valhalla, and will embody this possibility.


\flrightit{Posted on 2013-10-18 by Logres}

\begin{center}* * *\end{center}

\begin{footnotesize}
\begin{sffamily}
\texttt{William on 2013-10-19 at 12:06 said:}

``One of the greatest strengths and services and preparations for this is to read the exploits of the past ... ''

Can you suggest a resource or three? Thanx!

\hfill

\texttt{JA on 2013-10-19 at 12:45 said:}

I am ready.

\hfill

\texttt{Logres on 2013-10-19 at 18:04 said:}

William, part of the problem is ``teaching too soon, learning too late'' (Rosenstock-Huessy). This is a huge problem in Western modern education, generally speaking, and has lead to the current difficulty with ``elites'' who are only functionally so, in a shallow sense. Practically any good biography of a ``man of Tradition'' (or an auto-biography) from the last century or two will provide fodder for the young or old man who is attempting to advance. One thing it does it help us ``go back'' to the cross-roads and take the right path. Or, failing that, realize what was missed, and therefore, what extra labor must be undergone now to compensate for it. That is, there are some lessons of wisdom that have to be passed in the proper order to make spiritual advance, and studying the lives of the past helps fill in the blanks in modern psychology, deepening and widening and raising the standard. It also falls into the category of ``honoring your ancestors''. So we have to bring ``soul'' back into the endeavor. In the past, Plutarch's ``Lives'' was the standard place to start. Or, read someone's letters (such as Tolkien's or Coomaraswamy's, as Cologero has suggested). Evola's autobiography is ``Path of Cinnabar''. Tomberg's Meditations is a spiritual autobiography in a disguised form. Even when I was ignorant of Traditionalism, Martyn-Loyd Jones (a Welsh preacher) had a good biography from Banner of Truth, which helped me. You can also read Epictetus' maxims, or perhaps Aurelius' Meditations. There is no ``going wrong'' if you choose carefully and with enough mindfulness. The goal is to enter into their Mind, which is Logos (more or less). Of course, the biography of Christ is contained in the Gospels, and is argued to be also ``auto-biographical'' (God's account of his earthly Incarnation).

\hfill

\texttt{William on 2013-10-20 at 10:42 said:}

Thanks for the suggestions, Logres. I've already read about 1/3 of your suggestions. For what it's worth, I make part of my living as a composer, and lately I have been powerfully drawn to J.S.Bach --- not to imitate his style of composition or to have his music playing all day, but he was powerfully committed to the sacred craft of music --- more so than anyone else I know. What was his mindset? Etc. Trying to `get under his skin' a little has been very enlightening. Plutarch's Lives has been on my list of `I need to read this someday' but per your suggestion I'm thinking `someday' needs to be `order it now!' :-)

\hfill

\texttt{JA on 2013-10-20 at 13:41 said:}

I'd suggest, for practical purposes, if a man wants to be of the elite of the west, to really have the knowledge that an truly educated westerner did have in the past, he should begin by learning Latin. It's nice to read authors in translation but our ancestors only read Roman authors in the original. I'm working on Latin right now; Latin was the common language of Eternal pagan and Catholic Rome. afterwards he can try Greek... ... ... I'd say in reading, stay away from modern (post-1800 or so) literature, immerse yourself in the classics of Greece, Rome and the Middle Ages.

and... ... .at the same time we need to train our bodies. I'm trying to find or devise an entirely western based programme of martial arts for myself, based on chivalry. most of us cant do training on horseback but there should be enough resources to get our bodies fit for the fight.

\hfill

\texttt{Thorsten on 2013-10-20 at 14:12 said:}

JA,

If you haven't heard of them already ARMA (\url{http://www.youtube.com/user/theARMAonline}) has made great strides in reconstructing renaissance-era and late medieval martial arts. Most of it is limited to armed combat though, with lesser attention to unarmed Kampfringen (``combat grappling''). There are one or two Polish groups (E.G. Fechtschule Gdanks \url{http://www.youtube.com/watch?v=H7aXtzf7-Lk}) that appear to be pretty earnest. Russian Sambo and Systema would be more much more applicable to modernity, but as I understand it it is unclear how much of these martial arts are truly European in origin. However as the reconstructors of historical European martial arts have made clear there are only so many ways to disarm, throw, or kill a man so the most comprehensive fighting arts always end up closely resembling one another no matter where they originate.

\hfill

\texttt{Logres on 2013-10-20 at 20:20 said:}

William, don't leave out Nicholas Gomez Davila:

\url{http://don-colacho.blogspot.com/}

Easily scanned, as if Nietzsche had been a Christian. He has the rare talent of summarizing complex ideas in aphorisms, as if he was a Frenchman.

I haven't read Plutarch's Lives, but I have made a habit of repairing to antique texts as antidotes, and things are considerably safer when experimenting with a known quantity that the medievals and ancients also appreciated. Hard to see going wrong, on that one.

I would also suggest the study of falconry as something appropriate for the warrior class -- since that question was broached before about a year and a half ago. If you want something now, start planting seeds that will be NOW one day. Hector Paulus Mair is a big name in Western martial arts, in the ARMA revival. The quarterstaff, by the way, is an often overlooked ``martial arts'' weapon.

\hfill

\texttt{Logres on 2013-10-20 at 20:21 said:}

William, what (as a composer) do you think of Brahms?

\hfill

\texttt{kid on 2013-11-13 at 10:00 said:}

this series of articles is really excellent.

\hfill

\texttt{kid on 2013-11-20 at 10:42 said:}

there is western tradition found in any wrestling or boxing club. fencing too but thats probably out of the reach of most, myself included. besides these, judo and jiu jitsu are just jacketed wrestling with a little extra pomp and circumstance. sambo is pretty russian, but yah, borrowed a lot from judo. and it was created in the typical soviet way of doing things. that is, in a lab basically.

\hfill

\texttt{kid on 2013-11-20 at 10:44 said:}

savate as well!


\hfill

\end{sffamily}
\end{footnotesize}

